% !TeX root = ../main.tex
\section{Three generations of the same experiment}
\label{sec:generations}

The look-ahead comparison was run three times, in three independently built
implementations. No data flows between them: each generation simulates all of its own
conversations from scratch. Table~\ref{tab:generations} lists what changed.

\begin{table}[t]
\centering\small
\setlength{\tabcolsep}{4pt}
\begin{tabular}{@{}llll@{}}
\toprule
 & \textbf{Gen.\ 1} (ICLR'25) & \textbf{Gen.\ 2} & \textbf{Gen.\ 3} \\
\midrule
Therapist (policy)   & Llama-2-7B         & Llama-3.2-1B (4-bit NF4) & Llama-3.2-1B (bf16) \\
Patient simulator    & GPT-3.5            & \primaryjudge{}          & \primaryjudge{} \\
Patient prompts      & cooperative        & less cooperative         & less cooperative \\
Oracle (grader)      & GPT-3.5, regex     & \primaryjudge{}, JSON schema & \primaryjudge{}, JSON schema \\
Training reward      & mean(Q1, Q2)       & chosen oracle            & Q1{+}Q2 \\
Evaluation battery   & Q1, Q2             & 6 questionnaires         & 8 instruments \\
Held-out grader      & ---                & ---                      & \heldout{} \\
Optimizer            & PTO                & PTO                      & PTO \\
Look-ahead arms      & $\K \in \{0,5\}$   & $\K \in \{0,5\}$         & $\K \in \{0,5\}$ \\
Training oracles     & Q1{+}Q2            & Q1{+}Q2, WAI-SR, CSQ-8   & Q1{+}Q2 \\
Matched iterations   & 7                  & 5 (per oracle)           & 8 \\
Personas per model   & 96                 & 96                       & 96 \\
Persona order        & fixed              & fixed                    & reshuffled per iteration \\
\bottomrule
\end{tabular}
\caption{The three generations. Look-ahead depth $\K$ is the one lever measured identically
in all three, and all three use PTO --- which is what makes them a series. Generations~2 and~3
also ran other arms (a fourth control training oracle, an early static-data baseline, and in
generation~3 a full GRPO optimizer at $\K{=}0$); none of them contributes a matched $\K$
contrast, so none enters this chapter.}
\label{tab:generations}
\end{table}

\subsection{What each generation contributes}

\paragraph{Generation 1 (published).}
The ICLR experiment \citep{baruch2025pto}: a 7B therapist against GPT-3.5 patients, seven PTO
iterations at each of $\K{=}0$ and $\K{=}5$, scored on two questionnaires --- Q1 (session
satisfaction, 5 items) and Q2 (working alliance, 17 items), each on a 1--5 Likert scale, with
the training reward their mean. This is the generation whose look-ahead claim is under test.

\paragraph{Generation 2.}
A full re-implementation with a much smaller therapist, a stronger patient and oracle, less
cooperative patient personas, and a JSON-schema-constrained rubric covering six questionnaires.
Its distinctive contribution here is \emph{breadth}: it swept three training oracles
(Q1{+}Q2, WAI-SR, CSQ-8) at both look-ahead depths for five matched iterations each. Because
every model was scored on all six questionnaires, generation~2 alone supplies $105$
persona-paired $\K$ contrasts --- more than the other two generations combined.

One caveat matters for reading absolute numbers. Generation~2 generated its conversations in
4-bit NF4, which induces far more phrase-loop degeneration than the bf16 used in generation~3
(${\approx}9.5\%$ versus ${\approx}0.3\%$ of therapist turns running to the token cap as
repeated text). The oracle floors those conversations, so generation~2's scores sit roughly
$0.6$ points below generation~3's for what is otherwise the same base model. \textbf{Absolute
levels are therefore not comparable across generations 2 and 3.} The $\K$ contrast is
unaffected: both arms within generation~2 were generated under identical quantization, so the
paired difference is internally valid. This is the reason every claim in this chapter is a
\emph{within-generation} contrast rather than a cross-generation level comparison.

\paragraph{Generation 3.}
The most heavily instrumented of the three: bf16 generation, eight evaluation instruments (the
six questionnaires plus a patient change-talk coder and an MI-inconsistency coder),
persona-level pairing, measured oracle reproducibility, and a full re-scoring of every
conversation by \heldout{} --- a grader from a different model family that never played the
patient and never touched training. Its PTO look-ahead arm reaches eight matched iterations
against the myopic arm, the longest matched $\K$ comparison in this work.

\subsection{The one thing held constant}
Across all three generations the look-ahead manipulation itself is unchanged in kind: $\K$
alters only how much simulated future the oracle sees before scoring a candidate
(\S\ref{sec:pto}), and $\K{=}5$ appends five further utterances generated by the same policy
and patient that produced the prefix. Branch counts, filter thresholds and loss functions are
matched \emph{within} each generation's own $\K$ comparison. That is what licenses reading the
three as one series, and it is also the limit of what may be read: the generations differ in
several ways at once, so this chapter can establish \emph{that} the effect reverses and rule out
two specific explanations, but cannot attribute the reversal to any single factor
(\S\ref{sec:limitations}).
